\documentclass[10pt,a4paper]{amsart}
\usepackage[margin=19mm]{geometry}
\usepackage{amsmath,amssymb,mathtools}
\usepackage[scr=rsfs,cal=euler]{mathalfa}
\usepackage{xfrac}
\usepackage[unicode,hidelinks,bookmarks=false]{hyperref}
\newcommand{\ie}{i.e.\ }
\newcommand{\cf}{cf.\ }
\newcommand{\resp}{respectively\ }
\newcommand{\Subsection}{\subsection}
\providecommand{\nobreakdash}{\nobreak\hbox{-}}
\setlength{\emergencystretch}{2em}
\multlinegap0pt
\usepackage{amssymb}
\newtheorem{theorem}{Theorem}
\newcommand{\Inv}{\mathrm{inv}}
\newcommand{\Snk}[1]{S_{n,k}(#1)}
\newcommand{\Snkp}[1]{S'_{n,k}(#1)}
\title{Combinatorial study of the q-Catalan triangle and its generalizations}
\author{WIRDANE Youssouf}
\date{Source: arXiv:2605.14682v1 (CC BY 4.0)}
\begin{document}
\maketitle

\noindent\emph{Source and licence.} Combinatorial study of the q-Catalan triangle and its generalizations. Version arXiv:2605.14682v1 (CC BY 4.0), distributed by arXiv under the Creative Commons Attribution 4.0 International licence, http://creativecommons.org/licenses/by/4.0/, as recorded in the arXiv licence metadata for this exact version and shown on the abstract page https://arxiv.org/abs/2605.14682v1. Full licence notice: \texttt{licenses/arxiv-2605.14682-CC-BY-4.0.txt}.\par\smallskip

\noindent\textit{Editorial scope.} Counted unit: the article's theorem thm:qperm (source0.tex lines 263--269). The article's complete original proof of it is delivered with it (source0.tex lines 270--295, one \texttt{proof} environment of 26 lines). No mathematics is written, repaired or re-derived: the statement and the proof are the article's own lines, in the article's own notation, and no retained line is substituted or re-flowed.  No further source-local context is needed: every cross-reference of every retained line resolves inside the two delivered spans.  Nothing was omitted from any retained span, so the package carries no omission record.  No retained line contains a citation, so the wrapper carries no bibliography and no bibliography entry.  No retained line contains a figure environment, an image-inclusion command or drawing code, so no image is delivered and the package declares no asset dependency.  Editorial formatting only, in addition: an AMS \texttt{amsart} preamble that restates the article's own declarations and macros used by the retained lines, namely the environments \texttt{theorem} on separate counters, together with the standard packages those lines need.  The retained environments print in this document as Theorem 1 in order of appearance, whereas the article prints the same items with the numbering of its own full document.  The complete native source of the article as distributed in the arXiv e-print for this exact version is bundled unchanged as \texttt{sources/200605/source0.tex}.
\begin{theorem}[$q$-Catalan triangle via $312$-avoiding permutations]
\label{thm:qperm}
For all $n\ge1$ and $0\le k\le n$:
\begin{equation}\label{eq:inv312}
C_{n,k}(q) = \sum_{\pi\in\Snkp{312}} q^{\Inv(\pi)}.
\end{equation}
\end{theorem}

\begin{proof}
We proceed by induction on~$n$, using the partition
$\Snkp{312}=S'_{n,k-1}(312)\sqcup\Snk{312}$.

\textit{Base case.}
The set $S'_{n,0}(312)=S_{n,0}(312)$ consists of the unique
decreasing permutation $n(n{-}1)\cdots1$, which has $\Inv=\binom{n}{2}$;
hence $\sum_{\pi\in S'_{n,0}(312)}q^{\Inv(\pi)}=q^{\binom{n}{2}}=C_{n,0}(q)$.

\textit{Induction step.}
Define the insertion map
\[\Phi : S'_{n-1,k}(312)\to\Snk{312},\quad
\Phi(\pi_1\cdots\pi_{n-1})=\pi_1\cdots\pi_k\cdot n\cdot\pi_{k+1}\cdots\pi_{n-1}.\]
The map $\Phi$ is a bijection: $n$ preserves $312$-avoidance and fixes
the determining position at~$k$.
Since $n$ is larger than all other elements, its insertion at position $k+1$
creates exactly $n-k-1$ new inversions.
Therefore $\Inv(\Phi(\pi))=\Inv(\pi)+(n-k-1)$ for all $\pi\in S'_{n-1,k}(312)$.
By the induction hypothesis,
\[\sum_{\pi\in\Snk{312}} q^{\Inv(\pi)}
=q^{n-k-1}\sum_{\pi\in S'_{n-1,k}(312)} q^{\Inv(\pi)}
=q^{n-k-1}\,C_{n-1,k}(q).\]
The partition gives
$\sum_{\Snkp{312}}q^{\Inv}=C_{n,k-1}(q)+q^{n-k-1}C_{n-1,k}(q)=C_{n,k}(q)$,
which completes the induction.
\end{proof}

\end{document}
